\documentclass[11pt]{article}
\usepackage[T1]{fontenc}
\usepackage{lmodern,microtype,amsmath,amssymb,amsthm,mathtools,tikz-cd}
\usepackage[margin=1in]{geometry}
\usepackage{xurl,hyperref}
\hypersetup{colorlinks=true,linkcolor=blue,citecolor=blue,urlcolor=blue,
  pdftitle={Definability over a recursion category}}
\newtheorem{theorem}{Theorem}[section]
\newtheorem{corollary}[theorem]{Corollary}
\newtheorem{proposition}[theorem]{Proposition}
\theoremstyle{definition}
\newtheorem{definition}[theorem]{Definition}
\newcommand{\R}{\mathcal R}
\newcommand{\C}{\mathcal C}
\newcommand{\BDLat}{\mathbf{BDLat}}
\newcommand{\Tot}{\operatorname{Tot}}
\newcommand{\Comp}{\operatorname{Comp}}
\title{Definability over a recursion category}
\author{}
\date{29 October 2026}
\begin{document}
\maketitle

\begin{abstract}
In a nontrivial recursion category, representation of every complementary
domain pair by a total arrow implies that two distinguished predicates have
no separator given by a complemented domain. We formulate the representation
condition using an indexed bounded distributive lattice of predicates and a
natural monomorphism from the domain functor. The inseparability theorem
excludes disjoint domain extensions of the distinguished pair whose join is
the top predicate. In arithmetic, it implies that every consistent computably
enumerable extension of a consistent theory strongly defining every recursive
set numeralwise is incomplete.
\end{abstract}

\section{Computations and their domains}

We write $gf=g\circ f$, so that $f$ is applied first. Let $\R$ be a
locally small recursion category in Heller's prodominical sense: it has a
near-product structure and zero arrows, all its objects are isomorphic,
and it admits a Turing morphism \cite[Sections~5 and~9]{Heller}.
Write $\bar f:X\to X$ for the restriction of $f:X\to Y$.
The restriction is the domain of $f$, expressed as an endomorphism
\cite[Section~3]{DPH}; it satisfies the restriction-category axioms
\cite[Section~2.1.1]{RestrictionI}.
An arrow $f:X\to Y$ is \emph{total} when $\bar f=1_X$.

Write $X\boxtimes Y$ for the binary restriction product, with total
projections $\pi_1:X\boxtimes Y\to X$ and
$\pi_2:X\boxtimes Y\to Y$. For arrows $f:Z\to X$ and $g:Z\to Y$,
their pairing satisfies
\[
 \pi_1\langle f,g\rangle=f\bar g,
 \qquad \pi_2\langle f,g\rangle=g\bar f,
 \qquad \overline{\langle f,g\rangle}=\bar f\bar g.
\]
For $h:Z\to X\boxtimes Y$, the uniqueness equation is
$\langle\pi_1h,\pi_2h\rangle=h$.
Consequently the category $\C=\Tot(\R)$, with the same objects and only
the total arrows, has binary products. No terminal object in $\C$ is
assumed. Its diagonal is the total arrow
$\Delta_X=\langle1_X,1_X\rangle:X\to X\boxtimes X$.

A \emph{restriction idempotent}, or \emph{domain}, on $X$ is an
endomorphism $e:X\to X$ satisfying $\bar e=e$. Domains on $X$ are
ordered by $e\leq d$ if $de=e$, with meet $e\wedge d=ed$;
restriction idempotents commute. The zero and identity are the least
and greatest elements. Assume that the domains on each object form a
bounded distributive lattice and that, for every total $f:X\to Y$,
the assignment $e\mapsto\overline{ef}$ preserves finite joins.

Let $\BDLat$ denote the category of bounded distributive lattices and
homomorphisms preserving both bounds, finite meets, and finite joins.
The \emph{contravariant domain functor}
\[
 E:\C^{\mathrm{op}}\longrightarrow\BDLat
\]
is defined on objects by
\[
 E(X)=\{e:X\to X\mid\bar e=e\},
\]
with the lattice structure just specified. On a morphism
$f^{\mathrm{op}}:Y\to X$ of $\C^{\mathrm{op}}$, corresponding to a
total $f:X\to Y$ in $\C$, it is defined by
\[
 E(f^{\mathrm{op}})=f_E^*:E(Y)\longrightarrow E(X),
 \qquad f_E^*(e)=\overline{ef}.
\]
The restriction axioms give preservation of meets;
totality gives $f_E^*(1_Y)=1_X$, and restriction zeros give
$f_E^*(0_Y)=0_X$. Together with the finite-join hypothesis, these
identities make each $f_E^*$ a bounded lattice homomorphism.
For total $f:X\to Y$ and $g:Y\to Z$, the restriction axioms also give
$(1_X)_E^*=1_{E(X)}$ and $(gf)_E^*=f_E^*g_E^*$, the functor laws.
The restriction-idempotent functor and its associated fibration are
studied in \cite[Sections~4.1--4.2]{RestrictionI}.

A pair $(e_+,e_-)$ in $E(X)$ is \emph{complementary} when
\[
 e_+\wedge e_-=0_X,\qquad e_+\vee e_-=1_X.
\]
Write $\Comp(E(X))$ for the set of these pairs.

\section{An indexed lattice of predicates}

The domains need not include the predicates associated with a theory.
For example, the theorem set of a theory need not be computably enumerable.
We therefore supply another functor
\[
 P:\C^{\mathrm{op}}\longrightarrow\BDLat
\]
and a monomorphism
\[
 j:E\rightarrowtail P
 \quad\text{in }[\C^{\mathrm{op}},\BDLat].
\]
Equivalently, each component $j_X:E(X)\to P(X)$ is an injective
bounded lattice homomorphism. Each component is therefore an order
embedding: $j_X(e)\leq j_X(d)$ if and only if $e\leq d$.
An element of $P(X)$ is called a \emph{predicate on $X$}.
For a total $f:X\to Y$, write $f_P^*=P(f^{\mathrm{op}}):P(Y)\to P(X)$.
Naturality is the commutative square
\begin{equation}\label{eq:naturality}
\begin{tikzcd}[column sep=large]
 E(Y) \arrow[r,hook,"j_Y"] \arrow[d,"f_E^*"']
   & P(Y) \arrow[d,"f_P^*"] \\
 E(X) \arrow[r,hook,"j_X"'] & P(X).
\end{tikzcd}
\end{equation}
We omit the subscripts $E,P$ when the argument specifies the lattice.

The functor $P$ is an \emph{indexed bounded distributive lattice}. It
also has a description as a fibration with poset fibres. Define a category
$\int P$ whose objects are pairs $(X,p)$, where $p\in P(X)$, and whose
arrows $(X,p)\to(Y,q)$ are total arrows $f:X\to Y$ satisfying
$p\leq f^*q$. Projection onto $X$ defines a functor
$\int P\to\C$. For every $f:X\to Y$ and $q\in P(Y)$, its chosen
cartesian lift is
\[
 (X,f^*q)\xrightarrow{\ f\ }(Y,q).
\]
For $k:Z\to X$, an arrow $(Z,p)\to(Y,q)$ over $fk$ satisfies
$p\leq(fk)^*q=k^*f^*q$. This is precisely the condition for its
unique factorization over $k$ through the displayed lift.
Identity and composition of these lifts follow from functoriality of $P$.
The natural transformation $j$ induces a functor $\int E\to\int P$
over $\C$, by $(X,e)\mapsto(X,j_X(e))$, preserving the chosen lifts.

\section{Substitution and representation}

Fix objects $U,S$ of $\R$, a total arrow
$s:U\boxtimes U\to S$, and predicates $p_+,p_-\in P(S)$ satisfying
$p_+\wedge p_-=0$. In the arithmetic interpretation, $U$ supplies formula
and numeral codes, $S$ supplies sentence codes, $s$ performs numeral
substitution, and the distinguished predicates express provability and
refutability.

For a total arrow $g:U\to U$, put
\[
 c_g=s(g\boxtimes1_U):U\boxtimes U\longrightarrow S.
\]
\begin{definition}\label{def:rep}
A total arrow $g:U\to U$ \emph{represents} a complementary pair
$(e_+,e_-)\in\Comp(E(U))$ if
\begin{equation}\label{eq:representation}
 c_g^*p_+=\pi_2^*j_U(e_+),\qquad
 c_g^*p_-=\pi_2^*j_U(e_-)
 \quad\text{in }P(U\boxtimes U).
\end{equation}
The \emph{representation condition} requires that every complementary
pair in $E(U)$ have such a representative.
\end{definition}

\section{Inseparability}

A complementary pair $(d_+,d_-)\in\Comp(E(S))$ \emph{separates}
$p_+$ from $p_-$ when
\begin{equation}\label{eq:separator}
 p_+\leq j_S(d_+),\qquad p_-\leq j_S(d_-).
\end{equation}
\begin{theorem}\label{thm:inseparability}
Let $\R$ be a recursion category whose domain lattices have stable
finite joins as in Section~1, and put $\C=\Tot(\R)$.
Let $P:\C^{\mathrm{op}}\to\BDLat$ be a functor and
$j:E\rightarrowtail P$ a monomorphism in
$[\C^{\mathrm{op}},\BDLat]$, where $E$ is the domain functor.
Fix $U,S$, a total arrow
$s:U\boxtimes U\to S$, and disjoint $p_+,p_-\in P(S)$.
If $0_U\ne1_U$ and every complementary pair in $E(U)$ is represented
in the sense of Definition~\ref{def:rep}, then no complementary pair
in $E(S)$ separates $p_+$ from $p_-$.
\end{theorem}

\begin{proof}
Suppose that $(d_+,d_-)$ is a separator. Define the total arrow
$\delta=s\Delta_U:U\to S$. The pair
\[
 a_+=\delta^*d_-,\qquad a_-=\delta^*d_+
\]
is complementary, because total inverse image preserves the lattice
operations and bounds. Choose a total $g:U\to U$ representing this pair.
Put $\gamma=\langle1_U,g\rangle:U\to U\boxtimes U$ and
$t=\delta g:U\to S$. The product equations give $\pi_2\gamma=g$
and the commutative square
\[
\begin{tikzcd}[column sep=large,row sep=large]
 U \arrow[r,"\gamma"] \arrow[d,"g"']
   & U\boxtimes U \arrow[d,"c_g"] \\
 U \arrow[r,"\delta"'] & S.
\end{tikzcd}
\]
In fact $c_g\gamma=s\langle g,g\rangle=s\Delta_Ug=\delta g$.
Pulling \eqref{eq:representation} back along $\gamma$, and using
\eqref{eq:naturality}, yields
\[
 t^*p_+=g^*j_U(a_+)=j_U(t^*d_-),\qquad
 t^*p_-=g^*j_U(a_-)=j_U(t^*d_+).
\]
Pulling \eqref{eq:separator} back along $t$ now gives
\[
 j_U(t^*d_-)\leq j_U(t^*d_+),\qquad
 j_U(t^*d_+)\leq j_U(t^*d_-).
\]
Injectivity of $j_U$ implies $t^*d_+=t^*d_-$. These elements are
complementary because $t$ is total. Equal complementary elements have
both their meet and their join equal to the same element, so
$0_U=1_U$, contrary to the hypothesis.
\end{proof}

The theorem also holds for any category $\C$ with binary products,
functors $E,P:\C^{\mathrm{op}}\to\BDLat$, and a monomorphism
$j:E\rightarrowtail P$, provided that $E(U)$ is nontrivial and the
substitution and representation hypotheses hold. The same proof applies.

\section{Extensions and the numerical interpretation}

The theorem excludes complete disjoint extensions of the distinguished
pair whenever both new predicates arise from computational domains.
\begin{corollary}\label{cor:extension}
Under the hypotheses of Theorem~\ref{thm:inseparability}, suppose that
$b_+,b_-\in E(S)$ satisfy
\[
 p_+\leq j_S(b_+),\qquad p_-\leq j_S(b_-),\qquad
 j_S(b_+)\wedge j_S(b_-)=0.
\]
Then $j_S(b_+)\vee j_S(b_-)\ne1$.
\end{corollary}
\begin{proof}
If the join were $1$, preservation of the lattice operations and
injectivity of $j_S$ would imply
$b_+\wedge b_-=0_S$ and $b_+\vee b_-=1_S$.
The pair would be a separator, contrary to the theorem.
\end{proof}

For a numerical example, take the recursion category whose objects are
infinite recursive subsets of $\mathbb N$ and whose arrows are partial
recursive functions of the specified types, with extensional equality.
Fix a recursive bijection $\langle-,-\rangle:\mathbb N^2\to\mathbb N$
with recursive projections, and take
$X\boxtimes Y=\{\langle x,y\rangle:x\in X,\ y\in Y\}$.
The domains on $X$ are the computably enumerable subsets of $X$;
complementary domain pairs are recursive subsets of $X$ and their complements.
Take $P(X)$ to be the ordinary
powerset of $X$, with inverse image along total recursive functions,
and let $j_X$ include the domains as subsets.

Fix effective enumerations $a\mapsto\varphi_a(x)$ of unary arithmetic
formulas and $b\mapsto\sigma_b$ of sentences, together with a total
recursive function $s:\mathbb N^2\to\mathbb N$ satisfying
\[
 \sigma_{s(a,n)}=\varphi_a(\overline n).
\]
Here $\overline n$ is the numeral for $n$. Through the chosen pairing,
$s$ is a total arrow $U\boxtimes U\to S$, with $U=S=\mathbb N$.
For a consistent theory $T$, define
\[
 p_+=\{b:T\vdash\sigma_b\},\qquad
 p_-=\{b:T\vdash\neg\sigma_b\}.
\]
Consistency makes these predicates disjoint.

A formula $\varphi_a(x)$ \emph{strongly defines} a set
$A\subseteq\mathbb N$ \emph{numeralwise in $T$} when
\[
 n\in A\ \Longrightarrow\ T\vdash\varphi_a(\overline n),
 \qquad
 n\notin A\ \Longrightarrow\ T\vdash\neg\varphi_a(\overline n).
\]
For consistent $T$, these implications give the exact two preimage
equalities in \eqref{eq:representation}, with the constant recursive
map $g(x)=a$. Conversely, a representing $g$ yields a defining formula
with code $g(0)$. Thus representation of every complementary domain
pair is equivalent to strong numeralwise definability of every recursive
set in $T$.

An \emph{extension} of $T$ proves every theorem of $T$. A theory is
\emph{computably enumerable} if it admits a computably enumerable
axiomatization in the fixed effective proof calculus.

\begin{corollary}[Essential incompleteness]\label{cor:arithmetic}
Let $T$ be a consistent arithmetical theory in the effective language
above. If every recursive subset of $\mathbb N$ is strongly definable
numeralwise in $T$, then every consistent computably enumerable
extension of $T$ is incomplete.
\end{corollary}
\begin{proof}
The definability hypothesis supplies the representation condition.
Let $T'$ be a consistent computably enumerable extension of $T$.
Enumerating its proofs shows that its theorem and refutation sets are
computably enumerable domains
$b_+,b_-$ containing $p_+,p_-$, respectively. They are disjoint by
consistency. Corollary~\ref{cor:extension} says that their union is not
all of $\mathbb N$, which is precisely incompleteness of $T'$.
\end{proof}

\section{Morphisms and chosen operations}

Let primed symbols denote a second structure. Take a functor
$F:\R\to\R'$ preserving restriction, zero arrows, and the chosen
binary restriction products, and preserving finite joins of domains.
It restricts to $F_t:\C\to\C'$ and induces
$E(F)_X:E(X)\to E'(FX)$, with $E(F)_X(e)=F(e)$.
Require $FU=U'$, $FS=S'$, and $F(s)=s'$ under the chosen product
identifications. Supply a natural transformation
\[
 \eta:P\Longrightarrow P'\circ F_t^{\mathrm{op}}
\]
whose components are bounded lattice homomorphisms and satisfy
\[
 \eta_Xj_X=j'_{FX}E(F)_X,\qquad
 \eta_S(p_+)=p'_+,\qquad\eta_S(p_-)=p'_-.
\]
These equations ensure that $F(g)$ represents the image of every pair
represented by $g$. If Turing morphisms are part of the chosen recursion
data, require their preservation as well. Identity functors and
composition, with the corresponding predicate transformations, define
the category of these structures.

One may also specify a choice of representatives. Define the set
\[
 \operatorname{Rep}=
 \{(a,g)\in\Comp(E(U))\times\C(U,U):g\text{ represents }a\}.
\]
The representation condition says that the projection
$\operatorname{Rep}\to\Comp(E(U))$ is surjective. A chosen section
amounts to a function
\[
 \rho:\Comp(E(U))\longrightarrow\C(U,U)
\]
satisfying \eqref{eq:representation}. Ordinary set-theoretic choice
supplies such a section. In the numerical example, this is an external
choice of representing recursive functions; it does not assert a
recursive procedure for selecting their indices.

With fixed object, hom-set, and predicate sorts, the chosen structure has
the operations of composition, restriction, pairing, lattice operations,
inverse image, the maps $j_X$, and $\rho$. The domain of $\rho$ is specified
by the two complement equations. The one-object, domain-only version
takes $P=E$ and $j=1$.

A morphism preserving the chosen representatives satisfies
\[
 F_t(\rho(a))=\rho'(\Comp(E(F)_U)(a)).
\]
Thus choosing representatives adds a preservation equation to the
definition of a morphism.

\begin{thebibliography}{9}
\bibitem{DPH} R.~A.~Di Paola and A.~Heller,
Dominical categories: recursion theory without elements,
\emph{Journal of Symbolic Logic} \textbf{52}(3) (1987), 594--635.
\bibitem{Heller} A.~Heller,
An existence theorem for recursion categories,
\emph{Journal of Symbolic Logic} \textbf{55}(3) (1990), 1252--1268.
\bibitem{RestrictionI} J.~R.~B.~Cockett and S.~Lack,
Restriction categories I: categories of partial maps,
\emph{Theoretical Computer Science} \textbf{270}(1--2) (2002), 223--259.
\end{thebibliography}
\end{document}
